%-------------------------
% CV en LaTeX – Développeur full stack, architecture, DevOps et intégration IA
% Version finale optimisée pour Overleaf et les systèmes ATS
%-------------------------

\documentclass[a4paper,10pt]{article}

% Encodage français
\usepackage[T1]{fontenc}
\usepackage[utf8]{inputenc}

\usepackage{latexsym}
\usepackage[empty]{fullpage}
\usepackage{titlesec}
\usepackage{marvosym}
\usepackage[usenames,dvipsnames]{color}
\usepackage{verbatim}
\usepackage{enumitem}
\usepackage[hidelinks]{hyperref}
\usepackage{fancyhdr}
\IfFileExists{french.ldf}{
  \usepackage[french]{babel}
}{
  \usepackage[english]{babel}
}
\usepackage{tabularx}
\IfFileExists{fontawesome5.sty}{
  \usepackage{fontawesome5}
}{
  \newcommand{\faPhone}{}
  \newcommand{\faEnvelope}{}
  \newcommand{\faLinkedin}{}
  \newcommand{\faGithub}{}
}
\usepackage{multicol}
\usepackage{ragged2e}
\usepackage{geometry}

\geometry{
  a4paper,
  left=15mm,
  right=15mm,
  top=11mm,
  bottom=13mm
}

\setlength{\multicolsep}{-3.0pt}
\setlength{\columnsep}{-1pt}

\input{glyphtounicode}

\pagestyle{fancy}
\fancyhf{}
\fancyfoot{}
\renewcommand{\headrulewidth}{0pt}
\renewcommand{\footrulewidth}{0pt}

\urlstyle{same}
\raggedbottom
\raggedright
\setlength{\tabcolsep}{0in}
\setlength{\footskip}{10pt}
\setlength{\parindent}{0pt}
\setlength{\parskip}{0pt}
\setlength{\emergencystretch}{1em}
\linespread{1.0}\selectfont

\titleformat{\section}{\scshape\raggedright\large\bfseries}{}{0em}{}[\color{black}\titlerule]
\titlespacing*{\section}{0pt}{8pt}{5pt}

\pdfgentounicode=1

%-------------------------
% Commandes personnalisées
%-------------------------

\newcommand{\tagline}[1]{\small #1}

\newcommand{\resumeItem}[1]{
  \item\small{{#1 \vspace{-1pt}}}
}

\newcommand{\resumeSubheading}[4]{
  \vspace{0pt}\item
    \begin{tabular*}{1.0\textwidth}[t]{l@{\extracolsep{\fill}}r}
      \textbf{#1} & \textbf{\small #2} \\
      \textit{\small #3} & \textit{\small #4} \\
    \end{tabular*}\vspace{2.5pt}
}

\newcommand{\resumeProjectHeading}[2]{
  \item
    \begin{tabular*}{1.0\textwidth}{l@{\extracolsep{\fill}}r}
      \small #1 & \textbf{\small #2} \\
    \end{tabular*}\vspace{2.5pt}
}

\newcommand{\resumeSubItem}[2]{
  \item\small{\textbf{#1}: #2 \vspace{0.3pt}}
}

\renewcommand\labelitemi{$\vcenter{\hbox{\tiny$\bullet$}}$}
\renewcommand\labelitemii{$\vcenter{\hbox{\tiny$\bullet$}}$}

\newcommand{\resumeSubHeadingListStart}{
  \begin{itemize}[
    leftmargin=0.0in,
    label={},
    itemsep=1.5pt,
    topsep=2pt,
    parsep=0pt,
    partopsep=0pt
  ]
}

\newcommand{\resumeSubHeadingListEnd}{
  \end{itemize}
}

\newcommand{\resumeItemListStart}{
  \begin{itemize}[
    leftmargin=0.18in,
    itemsep=0.8pt,
    topsep=2pt,
    parsep=0pt,
    partopsep=0pt
  ]
}

\newcommand{\resumeItemListEnd}{
  \end{itemize}\vspace{1.5pt}
}

%-------------------------------------------
% Début du CV
%-------------------------------------------

\begin{document}


%----------EN-TÊTE----------
\begin{center}
    {\Huge \scshape Ibrahim Boubakri} \\ \vspace{1pt}
    \tagline{Développeur Full Stack \textbar{} DevOps \textbar{} Développement logiciel \textbar{} Intégration IA} \\ \vspace{2pt}
    \small
    \raisebox{-0.1\height}\faPhone\ +1 (819) 992-2001 ~
    \href{mailto:ibrahimboubakri1@gmail.com}{
    \raisebox{-0.2\height}\faEnvelope\ \underline{ibrahimboubakri1@gmail.com}} ~
    \href{https://linkedin.com/in/ibrahimboubakri/}{
    \raisebox{-0.2\height}\faLinkedin\ \underline{ibrahimboubakri}} ~
    \href{https://github.com/boubakriibrahim}{
    \raisebox{-0.2\height}\faGithub\ \underline{boubakriibrahim}}
    \vspace{-8pt}
\end{center}

%-----------PROFIL-----------
\section{Profil}

\justify
\small{
Développeur full stack orienté conception et livraison de solutions logicielles robustes,
performantes et maintenables, avec expérience en backend, frontend, APIs, automatisation,
DevOps et intégration IA. Conçoit des services modulaires, des workflows conteneurisés,
des pipelines CI/CD et des composants applicatifs enrichis par l'IA/LLM, de la conception
au déploiement et à l'observabilité. Solide maîtrise de Python, FastAPI, React, Node.js,
Docker, Kubernetes, SQL et Linux, avec un fort accent sur la fiabilité, l'efficacité,
la qualité logicielle, l'évolutivité et l'amélioration continue.
}

\vspace{-3pt}

%-----------FORMATION-----------
\section{Formation}

\resumeSubHeadingListStart

  \resumeSubheading
    {Université du Québec à Trois-Rivières -- UQTR}
    {Sept. 2024 -- Août 2026}
    {Maîtrise en génie électrique, concentration informatique}
    {Trois-Rivières, QC, Canada}

  \resumeSubheading
    {École Nationale des Sciences de l'Informatique -- ENSI}
    {Sept. 2020 -- Août 2023}
    {Diplôme d'ingénieur en informatique}
    {Université de la Manouba, Tunisie}

\resumeSubHeadingListEnd

%-----------EXPÉRIENCE PROFESSIONNELLE-----------
\section{Expérience professionnelle}

\resumeSubHeadingListStart

  \resumeSubheading
    {Développeur full stack -- Architecture, DevOps et intégration IA}
    {Sept. 2024 -- Aujourd'hui}
    {Modular AI Platform}
    {Québec, QC, Canada / Télétravail}

    \resumeItemListStart

      \resumeItem{
      Contribué à la conception et au développement d'une plateforme
      logicielle modulaire destinée à l'orchestration de traitements,
      à la gestion de données et à l'intégration de services
      d'intelligence artificielle.
      }

      \resumeItem{
      Conçu des composants backend et des APIs pour gérer les pipelines IA,
      les fichiers structurés, les métadonnées, les journaux d'exécution
      et les mécanismes de validation.
      }

      \resumeItem{
      Participé aux choix d'architecture afin de découpler les services,
      améliorer la maintenabilité du code et faciliter l'évolution
      indépendante des modules.
      }

      \resumeItem{
      Intégré des modèles et services d'IA dans des workflows applicatifs
      reproductibles, avec gestion des entrées, des sorties, des erreurs
      et du suivi d'exécution.
      }

      \resumeItem{
      Mis en place et amélioré des pratiques DevOps fondées sur la
      conteneurisation, l'intégration continue, l'automatisation des tests
      et la gestion des déploiements.
      }

      \resumeItem{
      Développé des mécanismes de stockage, de suivi d'avancement et de
      traitement asynchrone pour renforcer la fiabilité et l'observabilité
      de la plateforme.
      }

      \resumeItem{
      Collaboré à l'optimisation des performances, à la résolution
      d'incidents techniques et à la documentation des composants
      et interfaces.
      }

      \resumeItem{
      Accéléré le développement et la résolution d'incidents à l'aide
      d'assistants de programmation agentiques (Claude Code, OpenAI Codex,
      GitHub Copilot et Cursor) pour l'analyse de code, le refactoring,
      la génération de tests et la documentation, avec revue humaine
      et validation CI/CD.
      }

      \resumeItem{
      \textit{Technologies :}
      Python, FastAPI, REST, Docker, GitLab CI/CD,
      SQL, PyTorch, OpenCV, Linux
      }

    \resumeItemListEnd

    \vspace{7pt}



  \resumeSubheading
    {Chercheur -- Vision par ordinateur et robotique autonome}
    {2024 -- 2026}
    {UQTR \textbar{} Collaboration industrielle avec Noovelia}
    {Trois-Rivières, QC, Canada}

    \resumeItemListStart

      \resumeItem{
      Conçu et développé une chaîne de perception par vision de bout en bout
      pour la détection de palettes industrielles et l'estimation géométrique,
      avec un objectif de robustesse, d'efficacité et d'inférence temps réel.
      }

      \resumeItem{
      Développé un pipeline automatisé de données synthétiques sous Blender
      avec randomisation de domaine et annotations automatiques afin d'accélérer
      l'expérimentation et d'améliorer la généralisation vers des scènes réelles.
      }

      \resumeItem{
      Entraîné, ajusté et comparé des modèles YOLO-Pose avec PyTorch,
      Ultralytics et CUDA sur infrastructure GPU NVIDIA; structuré des protocoles
      reproductibles pour évaluer robustesse, généralisation et compromis performance/latence.
      }

      \resumeItem{
      Développé une plateforme intégrée reliant génération de données,
      entraînement, inférence, simulation PyBullet, services FastAPI/WebSocket
      et interface React afin d'unifier les workflows et réduire les opérations manuelles.
      }

      \resumeItem{
      Intégré le module de perception avec ROS 2 et une chaîne de navigation
      existante, déployé la solution sur NVIDIA Jetson, validé son fonctionnement
      en temps réel et préparé un pipeline réutilisable pour le partenaire industriel.
      }

      \resumeItem{
      \textit{Technologies :}
      Python, PyTorch, Ultralytics, CUDA, Blender, PyBullet, ROS 2,
      FastAPI, WebSocket, React, NVIDIA Jetson
      }

    \resumeItemListEnd

    \vspace{7pt}

  \resumeSubheading
    {Développeur full stack -- DevOps et automatisation}
    {Nov. 2023 -- Août 2024}
    {PRIMATEC Engineering}
    {Sfax, Tunisie}

    \resumeItemListStart

      \resumeItem{
      Développé des outils logiciels et des scripts d'automatisation
      pour réduire les opérations manuelles dans les processus de
      validation, de test et de contrôle qualité.
      }

      \resumeItem{
      Intégré des tests automatisés dans des pipelines CI/CD afin
      d'améliorer la fiabilité, la reproductibilité et la rapidité
      des livraisons logicielles.
      }

      \resumeItem{
      Conçu des composants Python modulaires avec gestion des erreurs,
      journalisation, traçabilité et séparation claire des responsabilités.
      }

      \resumeItem{
      Maintenu et amélioré des environnements d'intégration continue
      reposant sur Docker, Jenkins, GitLab et des scripts d'automatisation.
      }

      \resumeItem{
      Implémenté des mécanismes de collecte et d'exploitation de journaux
      techniques afin de faciliter le diagnostic des erreurs et le suivi
      des exécutions.
      }

      \resumeItem{
      Collaboré avec les équipes de développement et de validation
      pour standardiser les procédures techniques et sécuriser les
      étapes de livraison.
      }

      \resumeItem{
      Rédigé la documentation d'installation, d'utilisation et de
      maintenance des outils développés.
      }

      \resumeItem{
      \textit{Technologies :}
      Python, Docker, Jenkins, GitLab CI/CD, Jira, Batch,
      PowerShell, tests automatisés, Linux
      }

    \resumeItemListEnd

    \vspace{7pt}

  \resumeSubheading
    {Développeur full stack stagiaire -- Applications web et données}
    {Mars 2023 -- Août 2023}
    {Institut de Recherche sur l'Hydrogène}
    {Trois-Rivières, QC, Canada}

    \resumeItemListStart

      \resumeItem{
      Conçu et déployé une plateforme web reposant sur une architecture
      microservices pour la collecte, le traitement et la visualisation
      de données techniques en temps réel.
      }

      \resumeItem{
      Développé des services frontend et backend ainsi que des interfaces
      d'échange entre les composants de la plateforme.
      }

      \resumeItem{
      Implémenté des pipelines ETL automatisés pour collecter,
      transformer et stocker des flux de données transmis via MQTT.
      }

      \resumeItem{
      Conteneurisé les différents services afin d'améliorer la portabilité,
      l'isolation et la reproductibilité des environnements.
      }

      \resumeItem{
      Participé à l'intégration de pratiques CI/CD pour automatiser
      les phases de construction, de test et de déploiement.
      }

      \resumeItem{
      Contribué à la modélisation des données, à la gestion des bases
      de données et à la documentation de l'architecture technique.
      }

      \resumeItem{
      \textit{Technologies :}
      React, Node.js, MongoDB, Python, MQTT, Docker,
      Kubernetes, Jenkins
      }

    \resumeItemListEnd

\resumeSubHeadingListEnd

\vspace{6pt}

%-----------PROJETS-----------
\section{Projets sélectionnés}

\resumeSubHeadingListStart

  \resumeProjectHeading
    {\textbf{Plateforme modulaire IA et automatisation}}
    {Avr. 2026 -- Aujourd'hui}

    \resumeItemListStart

      \resumeItem{
      Conçu une architecture modulaire combinant services backend,
      APIs, automatisation et composants d'intelligence artificielle.
      }

      \resumeItem{
      Développé des workflows capables d'orchestrer plusieurs étapes
      de traitement, de validation et de génération de résultats.
      }

      \resumeItem{
      Implémenté des mécanismes de suivi des exécutions, de gestion
      des métadonnées, de versionnement des paramètres et de
      centralisation des sorties.
      }

      \resumeItem{
      Intégré des modèles d'IA générative dans des processus logiciels
      reproductibles avec validation des réponses et gestion des erreurs.
      }

      \resumeItem{
      Conteneurisé les services et automatisé les étapes de test
      et de déploiement afin de faciliter l'exploitation de la plateforme.
      }

      \resumeItem{
      \textit{Technologies :}
      Python, FastAPI, Docker, APIs REST, SQL,
      GitLab CI/CD, IA générative, LLM, RAG, MCP, tool/function calling
      }

    \resumeItemListEnd

    \vspace{7pt}

  \resumeProjectHeading
    {\textbf{HackZone Tunisia X}
    \textbar{} \emph{Infrastructure et cybersécurité}}
    {Avr. 2022}

    \resumeItemListStart

      \resumeItem{
      Participé à l'organisation d'un événement CTF de 24 heures
      réunissant plus de 100 équipes et 300 participants.
      }

      \resumeItem{
      Contribué à la conception et au déploiement de l'infrastructure
      technique utilisée pour héberger les défis et les services
      de la compétition.
      }

      \resumeItem{
      Préparé des défis portant sur les vulnérabilités web,
      les systèmes Linux et la rétro-ingénierie.
      }

      \resumeItem{
      Assuré le support technique, le suivi des services et la
      résolution d'incidents pendant l'événement.
      }

      \resumeItem{
      \textit{Technologies :}
      Azure, Python, Kali Linux, Git, Docker, cybersécurité web
      }

    \resumeItemListEnd

\resumeSubHeadingListEnd

\vspace{3pt}

%-----------COMPÉTENCES-----------
\section{Compétences techniques}

\resumeSubHeadingListStart

  \resumeSubItem{Développement full stack et backend}{
  Python, JavaScript, React, Node.js, FastAPI, Flask, APIs REST,
  WebSocket, services asynchrones, authentification, intégration de services
  et conception de composants réutilisables
  }

  \resumeSubItem{Architecture logicielle}{
  architectures modulaires, microservices, conception d'APIs,
  séparation des responsabilités, intégration de services,
  évolutivité, maintenabilité et optimisation des performances
  }

  \resumeSubItem{DevOps et livraison logicielle}{
  Git, GitLab CI/CD, Jenkins, Docker, Kubernetes, Linux,
  conteneurisation, tests automatisés, pipelines de déploiement,
  observabilité, journalisation et diagnostic d'incidents
  }

  \resumeSubItem{Intégration IA et LLM}{
  IA générative, intégration de LLM locaux et externes, RAG,
  MCP, tool/function calling, workflows agentiques, validation de réponses,
  PyTorch, OpenCV et automatisation assistée par IA
  }

  \resumeSubItem{Vision par ordinateur et robotique}{
  Ultralytics, YOLO-Pose, CUDA, ROS 2, NVIDIA Jetson,
  Blender, PyBullet, perception robotique, inférence temps réel
  et données synthétiques
  }

  \resumeSubItem{Données et pipelines}{
  SQL, MongoDB, MySQL, ETL, MQTT, JSON, CSV,
  validation de données, métadonnées, fichiers structurés
  et analyse de logs
  }

  \resumeSubItem{Programmation}{
  Python, SQL, JavaScript, Java, PHP, C++, C\#,
  PowerShell et Bash
  }

  \resumeSubItem{Outils de développement assisté par IA}{
  Claude Code, OpenAI Codex, GitHub Copilot, Cursor,
  analyse de code, refactoring, génération de tests,
  documentation et validation avec revue humaine
  }

  \resumeSubItem{Langues}{
  français, anglais, arabe
  }

\resumeSubHeadingListEnd

\vspace{-2pt}

%-----------CERTIFICATIONS ET ACTIVITÉS-----------
\section{Certifications et activités}

\resumeItemListStart

  \resumeItem{
  \textbf{Certifications :} DevOps Bootcamp -- TechWorld with Nana
  \textbar{} AWS Cloud Technical Essentials -- Coursera
  }

  \resumeItem{
  Membre de l'équipe d'organisation de ENSI Annual Forum,
  TuniHack, RoboCup et HackZone Tunisia.
  }

\resumeItemListEnd

\end{document}